\documentclass[12pt]{article}

\usepackage{bbm}
\usepackage{fancyhdr}
\usepackage{titling}
\usepackage{enumitem}
\usepackage{amsmath}
\usepackage{dsfont}
\usepackage{amsfonts}
\usepackage{amssymb}
\usepackage{amsthm}
\usepackage{geometry}
\usepackage{tikz}
\usepackage{algorithm}
\usepackage{algpseudocode}
\geometry{verbose,tmargin=1in,bmargin=1in,lmargin=1in,rmargin=1in}
\renewcommand{\vec}[1]{\mathbf{#1}}
\newcommand{\dec}{\text{dec}}
\newcommand{\divi}{\mathbin{\textnormal{div}}}
\title{Homework 1}
\author{Satvik Prasad}
\date{}
\newtheorem{lemma}{Lemma}[section]
\fancypagestyle{plain}{%
  \fancyhf{}
  \fancyhead[L]{\theauthor}
  \fancyhead[R]{\thedate}
}
\pagestyle{plain}
\makeatletter
\def\@maketitle{%
  \newpage
  \null
  \vskip 1em%
  \begin{center}%
    \let \footnote \thanks
    {\LARGE \@title \par}%
    \vskip 1em%
  \end{center}%
  \par
\vskip 1em}
\makeatother
\usepackage{cmbright}

\newcommand{\mat}[1]{\mathbf{#1}}

% ---------- Probability theory ----------
% Core operators (optional argument-free; use as \E{X}, \Var{X}, \Cov{X}{Y})
\DeclareMathOperator{\EOp}{\mathbb{E}}
\DeclareMathOperator{\VarOp}{Var}
\DeclareMathOperator{\CovOp}{Cov}
\DeclareMathOperator{\CorrOp}{Corr}

\renewcommand{\P}[1]{\mathbb{P}\left(#1\right)}        % \P{A}
\newcommand{\E}[1]{\EOp\left[#1\right]}                % \E{X}
\newcommand{\Var}[1]{\VarOp\left(#1\right)}            % \Var{X}
\newcommand{\Cov}[2]{\CovOp\left(#1,\,#2\right)}       % \Cov{X}{Y}
\newcommand{\Corr}[2]{\CorrOp\left(#1,\,#2\right)}     % \Corr{X}{Y}
\newcommand{\SD}[1]{\operatorname{SD}\left(#1\right)}  % \SD{X}

% Subscripted versions, e.g. \Esub{\theta}{X}, \Psub{\theta}{A}
\newcommand{\Esub}[2]{\EOp_{#1}\left[#2\right]}
\newcommand{\Psub}[2]{\mathbb{P}_{#1}\left(#2\right)}

% Conditioning, indicators, independence
\newcommand{\given}{\,\middle|\,}                      % \E{X \given Y}
\newcommand{\ind}[1]{\mathbbm{1}\left\{#1\right\}}     % \ind{X > 0}
\newcommand{\indep}{\perp\!\!\!\perp}                  % X \indep Y
\newcommand{\iid}{\overset{\text{iid}}{\sim}}          % X_i \iid \Normal(0,1)

% Convergence
\newcommand{\convp}{\xrightarrow{\;p\;}}               % in probability
\newcommand{\convd}{\xrightarrow{\;d\;}}               % in distribution
\newcommand{\convas}{\xrightarrow{\text{a.s.}}}        % almost surely
\newcommand{\convLp}[1]{\xrightarrow{L^{#1}}}          % in L^p

% Probability spaces
\newcommand{\F}{\mathcal{F}}                           % sigma-algebra
\newcommand{\B}{\mathcal{B}}                           % Borel sigma-algebra
\newcommand{\pspace}{(\Omega, \F, \mathbb{P})}

% Common distributions
\newcommand{\Normal}{\mathcal{N}}
\newcommand{\Bern}{\operatorname{Bernoulli}}
\newcommand{\Bin}{\operatorname{Binomial}}
\newcommand{\Pois}{\operatorname{Poisson}}
\newcommand{\Unif}{\operatorname{Uniform}}
\newcommand{\Exp}{\operatorname{Exp}}
\newcommand{\Geom}{\operatorname{Geometric}}
\newcommand{\Gam}{\operatorname{Gamma}}
\newcommand{\Beta}{\operatorname{Beta}}

% Number sets
\newcommand{\R}{\mathbb{R}}
\newcommand{\N}{\mathbb{N}}
\newcommand{\Z}{\mathbb{Z}}


\begin{document}
\begin{enumerate}
    \item Denote $A_{< j} = \bigcup_{i = 1}^{j -1}A_j.$
    Let $B_k = A_k \setminus A_{< k},$ and $B_1 = A_1.$ Then, we propose that 
    \begin{align*}
    \bigcup_{k = 1}^{n} A_k = \bigcup_{k = 1}^{n} B_k.
    \end{align*}

    This can be proved via double set containment. 
    
    Consider an element $a \in \bigcup_{k =1}^{n}A_k.$ Since $a$ is in this finite union, it is in at least one of $A_j$ for $1 \leq j \leq n.$ Take the smallest such $j.$ Then, $B_j$ contains $a$ since $a$ is not contained in $A_{< j},$ and is contained in $A_j.$ ($B_j = A_j \setminus A_{< j}$). Therefore, the finite union on the right hand side contains $a,$ since $B_j$ is in this union. Thus, $\bigcup_{k = 1}^{n}A_k \subseteq \bigcup_{k =1}^{n} B_k.$

    Conversely, consider an element $a \in \bigcup_{k =1 }^{n} B_k.$ Then $a \in B_j$ for some index $j.$ Because $B_j = A_j \setminus A_{<j},$ $a$ must be in $A_j,$ meaning $a$ is in $\bigcup_{k = 1}^{n} A_k.$ As a result, $\bigcup_{k =1}^{n}A_k \supseteq \bigcup_{k = 1}^{n} B_k,$ meaning the two sets are equal.

    Additionally, assume there is an element $b$ in both $B_i$ and $B_j$ for $i \neq j.$ WLOG, assume $i < j.$ Then, $b$ is in $A_i,$ but not in $A_{< i}.$ But that means, $b$ must be in $A_{< j}$ since $j > i,$ but that means it cannot be in $B_j.$ This a contradiction, meaning there is no element in $B_i \cap B_j$ and thus $B_i \cap B_j = \emptyset$ when $i \neq j.$

    Finally, since $(\Omega, \mathcal F, \mathbb P)$ is a probability measure, and $B_1, \ldots, B_n$ are pairwise disjoint, we can conclude \begin{align*}
        \mathbb P\left(
            \bigcup_{k =1 }^{n} A_k
        \right) = \mathbb P\left(
            \bigcup_{k = 1}^{n} B_k
        \right) = \sum_{k = 1}^{n} \mathbb P(B_k).
    \end{align*}

    But, notice that $A_k = B_k \cup A_{< k},$ meaning \begin{align*}
        \mathbb P(A_k) = \mathbb P(B_k \cup A_{< k}).
    \end{align*}
    Since $B_k$ and $A_{< k}$ are disjoint (by definition no element of $A_{<k}$ can be in $B_k$), we get
    \begin{align*}
        \mathbb P(A_k) = \mathbb P(B_k) + \mathbb P(A_{< k}) \geq \mathbb P(B_k),
    \end{align*}
    and thus we can conclude \begin{align*}
        \mathbb P\left(
            \bigcup_{k =1 }^{n} A_k
        \right) \leq \sum_{k = 1}^{n} \mathbb P(A_k)
    \end{align*}
    as required.

    \item 
    \begin{enumerate}
        \item When the outcome of every toss is of interest, any outcome can be represented via a binary string of length $10,$ meaning there are $2^{10}$ outcomes, and hence $2^{2^{10}}$ combinations of outcomes, meaning the size of our event space is $2^{2^{10}}.$ 

        \item In this case, an outcome can be represented via a single number ranging from $0$ to $10.$ Thus, there are $11$ outcomes, and thus our event space has size $2^{11}$ (power set of outcomes).
    \end{enumerate}
    
    \item \begin{enumerate}
        \item Probability $1 - (5/6)^4 \approx 0.518$
        \item Probability $1 - (35/36)^{24} \approx 0.491$
    \end{enumerate}

    Hence, (a) has a higher probability.

    \item Let $Y$ be the random variable representing the ordered selection of elements, and $X$ represent the unordered selection of elements. Notice that for $X = x,$ there are $3!$ selections for $Y = y$ that can achieve such an outcome. Hence, $$\mathbb P(X = x) = \sum_{y \text{ is a permutation of } x}\mathbb P(Y=  y).$$

    Now, we get $$
    \mathbb P(Y = y) = \frac{1}{9} \cdot \frac{1}{8} \cdot \frac{1}{7}, \qquad \mathbb P(X = x) = \dfrac{3!}{9 \cdot 8 \cdot 7} = \dfrac{1}{84}.
    $$

    There are 7 such unordered sets $x$ where all three numbers are adjacent (up to permutations), and hence $$\mathbb P(X \text{ adjacent}) = \dfrac{7}{84} = \frac{1}{12}.$$ 

    \item \begin{enumerate}
    \item The probability that no pairs of people share the same birthday is 
    $$
    \dfrac{365}{365} \cdot \ldots \cdot \dfrac{365 -k + 1}{365} = \dfrac{^{365}P_k}{365^{k}},
    $$
    meaning the probability that at least one pair of people share the same birthday is
    $$
    1- \dfrac{^{365}P_k}{365^{k}}.
    $$
    
    This is monotonically increasing. When $k = 23,$ the above evaluates to $\sim 0.507,$ whereas when $k = 22,$ the above evaluates to $\sim 0.476,$ meaning $k_\text{birthday} = 23.$

    \item Let $E_{ab}$ denote the event that people $a$ and $b$ share a birthday. The given argument computes
    $$
    \sum_{\{a, b\}} \mathbb P(E_{ab}) = {k \choose 2} \cdot \dfrac{1}{365},
    $$
    and treats this as $\mathbb P\left(\bigcup_{\{a,b\}} E_{ab}\right).$ But the events $E_{ab}$ are not mutually exclusive, so the probability of their union is not the sum of their probabilities.  E.g., $E_{ab}$ and $E_{bc}$ both occur whenever $a, b, c$ all share a birthday, and this outcome is counted twice in the sum. By the result of Question 1, the sum is only an upper bound:
    $$
    \mathbb P\left(\bigcup_{\{a,b\}} E_{ab}\right) \leq {k \choose 2} \cdot \dfrac{1}{365}.
    $$   
    \end{enumerate}

    \item 
    \begin{enumerate}
    \item We just expand twice: 
    \begin{align*}
    \mathbb P(A \cup B \cup C) &= \mathbb P(A \cup B) + \mathbb P(C) - \mathbb P((A \cup B) \cap C) \\
    &= \mathbb P(A) + \mathbb P(B) + \mathbb P(C) \\ &- \mathbb P(A \cap B) - \mathbb P((A \cap C) \cup (B \cap C)) \\
    &= \mathbb P(A) + \mathbb P(B) + \mathbb P(C) \\ &- \mathbb P(A \cap B) - \mathbb P(A \cap C) - \mathbb P(B \cap C) + \mathbb P(A \cap C \cap B \cap C) \\
    &= \mathbb P(A) + \mathbb P(B) + \mathbb P(C) \\ &- \mathbb P(A \cap B) - \mathbb P(A \cap C) - \mathbb P(B \cap C) \\ &+ \mathbb P(A \cap B \cap C)
    \end{align*}
    as required.

    \item Let $\sigma$ be a permutation of $(1, \ldots, n)$ sampled uniformly at random. We win if $\sigma(i) = i$ for some index $1 \leq i \leq n.$ Hence, we want to determine the probability that there is at least one such index.

    For each $i,$ let $$A_i = \{\sigma(i) = i\}$$ be the event that the $i$-th position is a fixed point. Then, our probability becomes
    $$
    p = \mathbb P(A_1 \cup A_2 \cup  \ldots \cup A_n).
    $$

    Notice that for any $\mathcal A \subseteq \{1, \ldots, n\},$ $$
    \mathbb P\left(\bigcap_{i \in \mathcal A} A_i\right) = \dfrac{(n - |\mathcal A|)!}{n!},
    $$ 
    since the positions in $\mathcal A$ are fixed and the remaining $n - |\mathcal A|$ positions may be permuted freely. Therefore, iterating over $|\mathcal A| = k$, we can use the principle of inclusion exclusion to yield
    $$
    p = \sum_{k = 1}^{n} (-1)^{k+1} \cdot \dfrac{(n - k)!}{n!} \cdot \underbrace{\left|\{\mathcal A \subseteq \{1, \ldots, n\} : |\mathcal A| = k\}\right|}_{n \choose k},
    $$
    simplifying down to 
    $$
    p = \sum_{k = 1}^{n} (-1)^{k+1} \dfrac{1}{k!}.
    $$

    For large $n,$ we get $$
    p \approx -\sum_{k = 1}^{\infty} \dfrac{(-1)^{k}}{k!} = -\left(e^{-1} - 1\right) = 1 - \dfrac{1}{e}.
    $$
    \end{enumerate}
\end{enumerate}
\end{document}